% ===== 原文 pp.23--28 =====
% [原文 p.23]
% （续 tr01b 末尾的同一列表：上链描述的第二至第四个要点）
\begin{itemize}[leftmargin=1.4em,itemsep=1pt,topsep=1pt]
\item $b_{d-p-1}$ 是一个 $\widehat{G}^{(p)}$ 取值的 $(d-p-1)$-上链，这意味着 $b_{d-p-1}$ 给三角剖分中的每个 $(d-p-1)$-单形指定阿贝尔群 $\widehat{G}^{(p)}$ 的一个元素。
\item $\delta$ 是一个微分，它把 $q$-上链映到取值于同一阿贝尔群的 $(q+1)$-上链
\begin{equation}
  \resizebox{0.88\textwidth}{!}{$\displaystyle \delta c_q(v_0,v_1,\cdots,v_{q+1}) = \sum_{i=0}^{q+1}(-1)^i c_q(v_0,v_1,\cdots,v_{i-1},v_{i+1},v_{i+2},\cdots,v_{q+1})$} \tag{2.84}
\end{equation}
其中左边是 $\delta c_q$ 在由顶点 $v_0,v_1,\cdots,v_{q+1}$ 所描述的 $(q+1)$-单形上的取值，右边是 $c_q$ 在 $(q+1)$-单形 $(v_0,v_1,\cdots,v_{q+1})$ 的各个子 $q$-单形上的取值。每个这样的子 $q$-单形都是通过从顶点 $v_0,v_1,\cdots,v_{q+1}$ 中删去一个顶点 $v_i$ 得到的。
\item 杯积 $\cup_\eta$ 一般取一个 $G$ 取值的 $q$-上链 $c_q$ 与一个 $H$ 取值的 $r$-上链 $d_r$，连同双同态
\begin{equation}
  \eta:\,G\times H\to K \tag{2.85}
\end{equation}
产生一个 $K$ 取值的 $(q+r)$-上链 $c_q\cup_\eta d_r$，它在 $(q+r)$-单形上的取值为
\begin{equation}
  \resizebox{0.88\textwidth}{!}{$\displaystyle c_q\cup_\eta d_r(v_0,v_1,\cdots,v_{q+r}) = \eta\left(c_q(v_0,v_1,\cdots,v_q)\in G,\ d_r(v_q,v_{q+1},\cdots,v_{q+r})\in H\right)\in K\,.$} \tag{2.86}
\end{equation}
\end{itemize}

用于作用量 (2.83) 中的杯积使用自然的双同态
\begin{equation}
  \eta:\,G^{(p)}\times\widehat{G}^{(p)}\to\R/\Z\,, \tag{2.87}
\end{equation}
其形式为
\begin{equation}
  \eta(g,\hat{g}) = \alpha\in\R/\Z\,,\qquad g\in G^{(p)},\ \hat{g}\in\widehat{G}^{(p)}\,, \tag{2.88}
\end{equation}
其中 $\alpha$ 通过
\begin{equation}
  \hat{g}(g) = \ee^{2\pi\ii\alpha}\in U(1) \tag{2.89}
\end{equation}
定义，这里用到了 $\hat{g}$ 本身是一个同态
\begin{equation}
  \hat{g}:\,G^{(p)}\to U(1) \tag{2.90}
\end{equation}
这一事实。

% [Tutor Note: p.23 (2.87)--(2.90) 是全部离散构造的关键：$\eta$ 就是 Pontryagin 配对。由 (2.89)，$\hat g(g)=\ee^{2\pi\ii\alpha}$ 只确定 $\alpha$ 模 $\Z$，所以 $\eta$ 的像确实落在 $\R/\Z$ 而不是 $\R$；这也正是 (2.83) 前面那个 $2\pi$ 因子的用处：作用量取 $2\pi\times(\R/\Z)$ 中的值，整体平移 $2\pi\Z$ 不改变 $\ee^{\ii S}$，故路径积分的权重良定义。可核对：把 $\hat g$ 换成 $\hat g^{\,n}$，则 $\eta$ 变为 $n\alpha$，仍是良定义的 $\R/\Z$ 元素。]

% [原文 p.24]
\begin{notebox}{标准杯积}
对
\begin{equation}
  G = H = K = \Z_n\,, \tag{2.91}
\end{equation}
人们常使用一个\textbf{标准杯积}，简记为 $\cup$，其形式为
\begin{equation}
  c_q\cup d_r(v_0,v_1,\cdots,v_{q+r}) = c_q(v_0,v_1,\cdots,v_q)d_r(v_q,v_{q+1},\cdots,v_{q+r})\,, \tag{2.92}
\end{equation}
其中右边我们使用了
\begin{equation}
  \Z_n = \Z/n\Z \tag{2.93}
\end{equation}
上的乘法，它由 $\Z$ 上通常的乘法下降而来。
\end{notebox}

% [Tutor Note: p.24 (2.92) 之所以要求 (2.91) 的 $G=H=K=\Z_n$，是因为只有 $\Z_n\times\Z_n\to\Z_n$ 的乘法才是双同态；对一般的 $G,H$ 并不存在这样的典范乘法，必须额外指定 $\eta$（这正是 (2.85) 的作用）。可核对：取 $G=H=\Z_2$、$K=\Z_4$，映射 $(a,b)\mapsto 2ab \bmod 4$ 才是双同态，而 $ab \bmod 4$ 不是。]

\para{规范变换。} 离散规范场 $a_p$ 是一个规范场，其规范变换作用为
\begin{equation}
  a_p\to a_p+\delta a_{p-1}\,, \tag{2.94}
\end{equation}
其中 $a_{p-1}$ 是任意一个 $G^{(p)}$ 取值的 $(p-1)$-上链。

类似地，离散规范场 $b_{d-p-1}$ 是一个规范场，其规范变换作用为
\begin{equation}
  b_{d-p-1}\to b_{d-p-1}+\delta b_{d-p-2}\,, \tag{2.95}
\end{equation}
其中 $b_{d-p-2}$ 是任意一个 $\widehat{G}^{(p)}$ 取值的 $(d-p-2)$-上链。

\para{运动方程。} 运动方程可以直接由作用量 (2.83) 算出：
\begin{equation}
  \delta a_p = 0\,,\qquad \delta b_{d-p-1} = 0\,, \tag{2.96}
\end{equation}
这意味着 $a_p$ 是一个 $p$-上闭链，而 $b_{d-p-1}$ 是一个 $(d-p-1)$-上闭链。

\para{高阶形式对称性。} 上面的运动方程使我们能够把 $a_p$ 与 $b_{d-p-1}$ 当作高阶形式 Noether 流的离散版本。也就是说，我们可以通过积分 $a_p$ 与 $b_{d-p-1}$ 来构造拓扑算符。例如，利用 $G^{(p)}$ 取值的 $p$-上闭链 $a_p$，我们可以构造由 $\widehat{G}^{(p)}$ 的元素标号的拓扑 $p$ 维算符
\begin{equation}
  U_{\hat{g}}(M_p) = \int_{M_p}\hat{g}(a_p)\,,\qquad \hat{g}\in\widehat{G}^{(p)}\,, \tag{2.97}
\end{equation}
其中 $\hat{g}(a_p)$ 是一个 $\U(1)$ 取值的 $p$-上链，它在 $p$-单形上的取值为
\begin{equation}
  (\hat{g}(a_p))(v_0,v_1,\cdots,v_p) = \hat{g}\left(a_p(v_0,v_1,\cdots,v_p)\in G^{(p)}\right)\in U(1)\,. \tag{2.98}
\end{equation}

% [原文 p.25]
上述积分于是被表示成
\begin{equation}
  \int_{M_p}\hat{g}(a_p) = \prod_{S_p\in M_p}(\hat{g}(a_p))(S_p)\,, \tag{2.99}
\end{equation}
其中 $S_p$ 是子流形 $M_p$ 的三角剖分中的单形，而 $(\hat{g}(a_p))(S_p)$ 是 $\U(1)$ 取值的 $p$-上链 $\hat{g}(a_p)$ 在 $S_p$ 上的取值。利用 Stokes 定理的离散版本以及运动方程 $\delta a_p=0$，可以看出这些算符 $U_{\hat{g}}$ 是拓扑的。这些算符 $U_{\hat{g}}$ 有时被称作规范场 $a_p$ 的 \textbf{Wilson 算符}。

类似地，利用 $\widehat{G}^{(p)}$ 取值的 $(d-p-1)$-上闭链 $b_{d-p-1}$，我们可以构造由 $G^{(p)}$ 的元素标号的拓扑 $(d-p-1)$ 维算符
\begin{equation}
  U_g(M_{d-p-1}) = \int_{M_{d-p-1}} b_{d-p-1}(g)\,,\qquad g\in G^{(p)}\,, \tag{2.100}
\end{equation}
其中 $b_{d-p-1}(g)$ 是一个 $\U(1)$ 取值的 $(d-p-1)$-上链，它在 $(d-p-1)$-单形上的取值为
\begin{equation}
  \resizebox{0.88\textwidth}{!}{$\displaystyle (b_{d-p-1}(g))(v_0,v_1,\cdots,v_{d-p-1}) = \left(b_{d-p-1}(v_0,v_1,\cdots,v_{d-p-1})\in\widehat{G}^{(p)}\right)(g)\in \U(1)\,.$} \tag{2.101}
\end{equation}
作为运动方程 $\delta b_{d-p-1}=0$ 的推论，这些算符是拓扑的。这些算符 $U_g$ 有时被称作规范场 $b_{d-p-1}$ 的 \textbf{Wilson 算符}。

总的来说，离散高阶形式规范理论 (2.83) 具有两类高阶形式对称性：
\begin{itemize}[leftmargin=1.4em,itemsep=1pt,topsep=1pt]
\item 一个 $p$-形式对称群
\begin{equation}
  G^{(p)}\,, \tag{2.102}
\end{equation}
它由余维-$(p+1)$ 拓扑算符 (2.100) 构成。
\item 一个 $(d-p-1)$-形式对称群
\begin{equation}
  G^{(d-p-1)} = \widehat{G}^{(p)}\,, \tag{2.103}
\end{equation}
它由 $p$ 维拓扑算符 (2.97) 构成。
\end{itemize}
这两类高阶形式对称性有时也称为电与磁对称性，以示区分。

\para{带电算符。} 一个 $p$ 维拓扑算符 $U_{\hat{g}}$ 携带荷
\begin{equation}
  Q(U_{\hat{g}}) = \hat{g}\in\widehat{G}^{(p)} \tag{2.104}
\end{equation}
% [原文 p.26]
位于由拓扑算符 $U_g$ 生成的 $G^{(p)}$ $p$-形式对称性之下。这可以看作下述事实的推论：$b_{d-p-1}$ 的运动方程在存在 $U_{\hat{g}}(M_p)$ 时被修改为
\begin{equation}
  U_{\hat{g}}(M_p)\delta b_{d-p-1} = \delta^{d-p}_{\hat{g}}(M_p)U_{\hat{g}}(M_p)\,, \tag{2.105}
\end{equation}
其中 $\delta^{d-p}_{\hat{g}}(M_p)$ 是局域在 $M_p$ 上的 delta 函数的一个离散类比物。它是一个 $\widehat{G}^{(p)}$ 取值的 $(d-p)$-上链，具有如下性质：
\begin{equation}
  \int\delta^{d-p}_{\hat{g}}(M_p)\cup_\eta\Lambda_p = \int_{M_p}\hat{g}(\Lambda_p)\in\R/\Z \tag{2.106}
\end{equation}
对任意 $G^{(p)}$ 取值的 $p$-上链 $\Lambda_p$ 成立。

% [Tutor Note: p.26 (2.106) 里杯积的次序是 $\delta^{d-p}_{\hat g}\cup_\eta\Lambda_p$，即 $\widehat{G}^{(p)}$ 取值在前、$G^{(p)}$ 取值在后，而 (2.87)/(2.88) 定义的 $\eta$ 是 $G^{(p)}\times\widehat{G}^{(p)}\to\R/\Z$。两者差一次对调；严格说应写 $\eta(\Lambda_p,\delta^{d-p}_{\hat g})$，或声明 $\eta$ 的变元次序按上下文理解。(2.109) 亦然。这不影响等式成立：只要令 $\eta'(\hat g,g):=\eta(g,\hat g)$ 即可，数值完全一样，属于原文记号次序不严。]

类似地，一个 $(d-p-1)$ 维拓扑算符 $U_g$ 携带荷
\begin{equation}
  Q(U_g) = g\in\widehat{G}^{(d-p-1)} = G^{(p)} \tag{2.107}
\end{equation}
位于由拓扑算符 $U_{\hat{g}}$ 生成的 $G^{(d-p-1)}=\widehat{G}^{(p)}$ $(d-p-1)$-形式对称性之下。这可以看作下述事实的推论：$a_p$ 的运动方程在存在 $U_g(M_{d-p-1})$ 时被修改为
\begin{equation}
  U_g(M_{d-p-1})\delta a_p = \delta^{p+1}_{g}(M_{d-p-1})U_g(M_{d-p-1})\,, \tag{2.108}
\end{equation}
其中 $\delta^{p+1}_{g}(M_{d-p-1})$ 是一个 $G^{(p)}$ 取值的 $(p+1)$-上链，具有性质
\begin{equation}
  \resizebox{0.88\textwidth}{!}{$\displaystyle \int\Lambda_{d-p-1}\cup_\eta\delta^{p+1}_{g}(M_{d-p-1}) = \int_{M_{d-p-1}}\Lambda_{d-p-1}(g)\in\R/\Z\,,$} \tag{2.109}
\end{equation}
对任意 $\widehat{G}^{(p)}$ 取值的 $(d-p-1)$-上链 $\Lambda_{d-p-1}$ 成立。

% [Tutor Note: p.26 (2.103) 与 (2.107) 自洽性核对：由 (2.103) $G^{(d-p-1)}=\widehat{G}^{(p)}$，再取一次对偶并用双重 Pontryagin 对偶 (2.74) 得 $\widehat{G}^{(d-p-1)}=\widehat{\widehat{G}}^{(p)}=G^{(p)}$，正是 (2.107) 的最后一步。若把 (2.107) 写成 $\widehat{G}^{(d-p-1)}=G^{(p)}$ 而不能反过来，就会与 (2.74) 冲突——这是检查此处上下标是否抄错的最快办法。]

\section{规范理论的高阶形式对称性}

\subsection{延展算符的屏蔽}

在本小节中，我们讨论一种理解高阶形式对称性的有力而简单的方法；在随后各小节中，我们将用它来计算规范理论的高阶形式对称性。这一方法以“屏蔽”（screening）之名著称，近来已被成功地用于各种强耦合情形，以推断高阶形式对称性 [62--76]。

我们从屏蔽的定义开始：

\begin{definitionbox}{定义 3.1：屏蔽}
我们说一个 $p\ge1$ 维算符 $\cO_p$ 可以被\textbf{屏蔽}（screened）成另一个 $p$ 维算符 $\cO_p'$，如果可以在 $\cO_p$ 与 $\cO_p'$ 之间插入一个 $(p-1)$ 维算符 $\cO_{p-1}$。见图 7。

如果算符 $\cO_p$ 可以被屏蔽成恒等 $p$ 维算符，即如果存在一个可以用来终止 $\cO_p$ 的算符 $\cO_{p-1}$，那么我们就简单地说 $\cO_p$ 可以被\textbf{完全屏蔽}。
\end{definitionbox}

% [原文 p.27]
\begin{figure}[htbp]
\centering
\includegraphics[width=0.68\textwidth]{fig7.png}
\figtitle{7}{上述构型的存在意味着 $\cO_p$ 与 $\cO_p'$ 由屏蔽相联系。}
\end{figure}

\para{屏蔽意味着相同的 $p$-形式荷。} 上述定义之所以有意义，是因为：若 $\cO_p$ 可被屏蔽成 $\cO_p'$，则这两个算符在任意 $p$-形式对称群 $G^{(p)}$ 下必须具有相同的荷
\begin{equation}
  Q(\cO_p) = Q(\cO_p')\in\widehat{G}^{(p)}\,. \tag{3.1}
\end{equation}
这一点在图 8 中得到说明。若一个算符 $\cO_p$ 可被完全屏蔽，则它在 $G^{(p)}$ 下不带荷
\begin{equation}
  Q(\cO_p) = 1\in\widehat{G}^{(p)}\,. \tag{3.2}
\end{equation}

% [Tutor Note: p.27 (3.1)--(3.2) 只给出单向推论（可屏蔽 ⇒ 荷相同；完全屏蔽 ⇒ 无荷），反向不成立：荷相同并不保证两算符由屏蔽相联系。可核对：在纯 $\U(1)$ 规范理论中任何两条同电荷 Wilson 线的 1-形式荷相同，但它们一般并不同类；故 (3.2) 只能用来“排除”某个算符带荷，不能用来“证明”它可被完全屏蔽。]

\para{由无屏蔽得到 $p$-形式对称性。} 反过来，若 $\cO_p$ 与 $\cO_p'$ 不是由屏蔽相联系的，那么我们就可以定义一个 $p$-形式对称性，其刻画性质正是：$\cO_p$ 与 $\cO_p'$ 在这一 $p$-形式对称性下携带不同的荷。

更具体地说，可以采用如下策略。在 $p$ 维算符的集合上定义一个等价关系
\begin{equation}
  \cO_p\sim\cO_p' \tag{3.3}
\end{equation}
如果 $\cO_p$ 与 $\cO_p'$ 由屏蔽相联系。令
\begin{equation}
  \mathsf{D}_p \tag{3.4}
\end{equation}
为施加上述等价关系后所得的 $p$ 维算符等价类的集合。这一集合从 $p$ 维算符的 OPE 获得一个乘积结构，从而被转换成一个环。在许多感兴趣的情形中，乘积结构受到更多限制，$\mathsf{D}_p$ 实际上是一个阿贝尔群。以下我们假定属于这种情形。

由上述讨论，我们必然有 $p$-形式对称性，其荷能够区分 $\mathsf{D}_p$ 的不同元素。这样的 $p$-形式对称性构成群
\begin{equation}
  G^{(p)} = \widehat{\mathsf{D}}_p\,. \tag{3.5}
\end{equation}
在这一 $p$-形式对称群下可能的荷是
\begin{equation}
  \widehat{G}^{(p)} = \mathsf{D}_p \tag{3.6}
\end{equation}
的元素。

% [Tutor Note: p.27 (3.4)--(3.6) 中 $\mathsf{D}_p$ 用的是无衬线的正体 D（不是 $\mathcal{D}$、也不是斜体 $D$），与原文一致；$\widehat{\mathsf{D}}_p$ 的帽子只盖住 D、下标 $p$ 在帽外，位置也照原文。自洽性核对：(3.5) 与 (3.6) 互为对偶，把 (3.5) 再取一次对偶即得 $\widehat{G}^{(p)}=\widehat{\widehat{\mathsf{D}}}_p=\mathsf{D}_p$，这正是 (3.6)，用到了 (2.74)。]

% [原文 p.28]
\begin{figure}[htbp]
\centering
\includegraphics[width=0.52\textwidth]{fig8.png}
\figtitle{8}{上方是两个仅由拓扑形变 $p$-形式对称性算符 $U_g$ 而相互联系的构型。下方是我们分别把 $U_g$ 收缩到 $\cO_p$ 与 $\cO_p'$ 上，得到相位 $Q(\cO_p)(g)$ 与 $Q(\cO_p')(g)$，其中 $Q(\cO_p),Q(\cO_p')\in\widehat{G}^{(p)}$。由于 $g\in G^{(p)}$ 是任取的，这便意味着 (3.1)。}
\end{figure}

% [Tutor Note: p.28 图 8 证明 (3.1) 的论证只用了一条输入：$U_g$ 是拓扑算符，故图 8 上方两构型相等（即 (2.39)）。收缩得到的两相位因此对同一个 $g$ 必须相等；又 $g$ 任意，故 $Q(\cO_p)=Q(\cO_p')$ 作为 $\widehat{G}^{(p)}$ 的元素相等。注意这里用到 $\widehat{G}^{(p)}$ 是阿贝尔群，从而“荷相等”与“相位相等”等价——这正是 $p\ge1$ 时 $G^{(p)}$ 必为阿贝尔（图 4）的用处。]
